% \iffalse meta-comment
%%
%% semantic-markup.dtx
%%
%% Copyright (C) 2026 Seymour J. Metz
%%
%% This material is subject to the MIT License.
%% See the LICENSE file in the root of this distribution, or
%% <https://github.com/shmuelmetz/LaTeX-Semantic-Markup>.
%%
%% This is the semantic-markup package, part of the
%% LaTeX-Semantic-Markup distribution.
%%
%<*driver>
\documentclass[full]{l3doc}
\usepackage{mathrsfs}
\usepackage{bm}
\usepackage{amssymb}
\begin{document}
  \DocInput{semantic-markup.dtx}
\end{document}
%</driver>
% \fi
%
% \title{The \pkg{semantic-markup} package}
% \author{Shmuel (Seymour J.) Metz \\ \url{https://mason.gmu.edu/~smetz3}}
% \date{2026-08-27}
%
% \maketitle
%
% \begin{documentation}
%
% \section{Purpose}
%
% This package provides macros that record the \emph{semantic role}
% of a piece of mathematical text --- ``this is the display name of a
% category,'' ``this is a monomorphism arrow'' --- separately from
% how that role is \emph{typeset}. An author writing
% \cs{catName}\Arg{C} is not choosing a font; they are asserting that
% \texttt{C} names a category, and the package decides (by default,
% and later by configuration) what a category name looks like.
%
% This separation is the entire point of the package: two things
% that look similar on the page but mean different things
% mathematically must not share a macro, even if that means near-
% identical default typesetting for two commands. Conversely, one
% concept must not be given two names merely because it is
% occasionally styled differently in context.
%
% \section{Internal naming}
%
% Package-internal \pkg{expl3} functions use the module prefix
% \texttt{lsm} (for \emph{latex-semantic-markup}), per the naming
% convention \texttt{\textbackslash lsm\_}\emph{module}\texttt{\_}%
% \emph{description}\texttt{:}\emph{args}. These are not intended for
% direct use in documents.
%
% \section{Abbreviations used in macro names}
%
% Every abbreviated word appearing in a user-level macro name is
% spelled out here, once, so no reader has to guess it from context.
% Family and role prefixes that are already complete English words
% (\texttt{chart}, \texttt{atlas}, \texttt{name}, \texttt{compose})
% are not abbreviations and are not listed.
%
% \begin{description}
% \item[\texttt{sem}] Semantic
% \item[\texttt{cat}] Category
% \item[\texttt{obj}] Object
% \item[\texttt{morph}] Morphism
% \item[\texttt{coord}] Coordinate
% \item[\texttt{seq}] Sequence (as in \cs{catSeqName}: a
%   \emph{sequence of} the preceding concept, not that concept itself)
% \item[\texttt{mono}] Monomorphism
% \item[\texttt{epi}] Epimorphism
% \end{description}
%
% \section{User-level command families}
%
% User-level commands are grouped by the semantic domain they tag,
% each using its own name prefix:
% \begin{itemize}
%   \item \texttt{sem} (Semantic) --- generic / not-yet-classified
%     semantic text
%   \item \texttt{cat} (Category) --- categories (in the
%     category-theoretic sense)
%   \item \texttt{obj} (Object) --- objects of a category
%   \item \texttt{morph} (Morphism) --- morphisms between objects
%   \item \texttt{chart} --- local coordinate charts
%   \item \texttt{atlas} --- atlases of coordinate charts
%   \item \texttt{coord} (Coordinate) --- coordinate systems/functions
% \end{itemize}
%
% \subsection{Generic}
%
% \DescribeMacro{\semName}
% \cs{semName}\Arg{text} (Sem = Semantic) tags \texttt{text} as
% semantically significant without committing to any of the more
% specific domains below. Use this only as a placeholder until a
% more specific macro applies, or for genuinely domain-agnostic
% labels.
%
% \subsection{Categories}
%
% \DescribeMacro{\catName}
% \cs{catName}\Arg{text} (Cat = Category) is the display name of a
% single category. The default presentation (\cs{mathscr}) follows
% the source papers' own stated convention: ``Upper case Script Latin
% letters will refer to categories and functors.''
%
% \DescribeMacro{\catSeqName}
% \cs{catSeqName}\Arg{text} (Cat = Category, Seq = Sequence) is the
% display name of a \emph{sequence or family of categories}, e.g.
% $(\mathscr{A}_\alpha)_{\alpha \in A}$. This is deliberately a
% distinct command from \cs{catName}, not a styling variant of it: a
% sequence of categories is not itself a category, and collapsing the
% two into one macro differentiated only by an optional ``bold''
% switch would silently re-introduce the typography/semantics
% conflation this package exists to remove. The source papers draw
% the same line themselves, independently, in more than one macro
% family: their own \texttt{catname}/\texttt{catseqname} pair is
% introduced under the single comment ``Select font for categories
% and sequences of categories,'' and separately their
% \texttt{subcat}/\texttt{SUBCAT} pair is commented ``Subcategory'' /
% ``Sequence of subcategories.'' The single-vs-sequence distinction
% is a recurring pattern in the source mathematics, not a one-off.
%
% \begin{quote}
% \textbf{Open question, not yet resolved:} the source papers also
% use a separate, bold-roman \texttt{cat} macro (distinct from
% \texttt{catname}) specifically for \emph{standard, well-known}
% categories --- their own comments say ``Use \cs{cat}\Arg{Top} for
% the standard category Top'' and likewise for \cs{cat}\Arg{Set}.
% That is a third role this package does not yet provide for: naming
% a specific, conventionally-fixed category (Set, Top, \ldots) is
% arguably semantically different from naming an arbitrary/variable
% category symbol via \cs{catName}, in the same way that
% \cs{catSeqName} is different from \cs{catName}. Whether that
% deserves its own macro (e.g.\ \cs{catStandardName}) or is genuinely
% just a presentational choice (bold vs.\ script) is left to the
% author to decide; it is not built here.
% \end{quote}
%
% \subsection{Objects}
%
% \DescribeMacro{\objName}
% \cs{objName}\Arg{text} (Obj = Object) is the display name of an
% object of a category, corresponding to what the source papers call
% $\Ob(\catname{C})$ --- ``object of category'' by their own comment,
% ``all objects of $\catname{C}$'' by their own formal definition.
% The default presentation is unstyled (objects are typically bare
% math-mode variables in practice); the macro exists so that document
% text records the semantic role even when the default rendering does
% nothing visible yet.
%
% \subsection{Morphisms}
%
% \DescribeMacro{\morphName}
% \cs{morphName}\Arg{text} (Morph = Morphism) is the display
% name/label of a morphism (e.g.\ a function symbol). The
% \texttt{morph} prefix follows the source papers' own explicit
% terminology choice: ``This paper uses the term morphism in
% preference to arrow, but uses the conventional $\Ar$.''
%
% \begin{quote}
% \textbf{Open question, not yet resolved:} the source papers define
% four related but distinct morphism-level operators that this
% package does not yet provide: $\Ar(\catname{C})$ and
% $\Mor(\catname{C})$ (``all morphisms of $\catname{C}$,'' explicitly
% stated as synonyms of each other), and
% $\Hom_\catname{C}(a,b)$ (``all morphisms of $\catname{C}$ from $a$
% to $b$'' --- morphisms between two \emph{specific} objects, a
% narrower notion than $\Ar$/$\Mor$). These are real, well-defined,
% distinct mathematical notions in the source and are natural
% candidates for something like \cs{morphAll} and \cs{morphBetween},
% but are not built here pending the author's naming preference.
% Note also a discrepancy in the source itself: the inline comment
% directly above the papers' \texttt{Mor} definition reads
% ``Morphisms between two objects'' (i.e.\ the \emph{Hom} reading),
% while the papers' own later formal prose defines \texttt{Mor} as
% ``all morphisms of $\catname{C}$, synonymous with $\Ar$'' (i.e.\
% \emph{not} the Hom reading). This looks like a stray/copy-pasted
% comment in the source rather than a real ambiguity in the
% mathematics, but it is flagged here rather than silently resolved
% either way.
% \end{quote}
%
% \DescribeMacro{\morphMono}
% \cs{morphMono} (Mono = Monomorphism) is the arrow symbol for a
% monomorphism, corresponding to the source papers' \texttt{into}.
%
% \DescribeMacro{\morphEpi}
% \cs{morphEpi} (Epi = Epimorphism) is the arrow symbol for an
% epimorphism, corresponding to the source papers' \texttt{onto}.
% This is kept distinct from \cs{morphMono} because mono/epi is a
% real mathematical distinction (injective vs.\ surjective in the
% relevant sense), not a choice of arrowhead glyph.
%
% \DescribeMacro{\morphCompose}
% \cs{morphCompose} is the general morphism-composition operator
% (default $\circ$). An optional argument
% \cs{morphCompose}\oarg{label} places \texttt{label} over the
% operator, mirroring how composition is annotated in the source
% papers.
%
% \begin{quote}
% \textbf{Open question, not yet resolved:} the source papers define
% \emph{three} composition operators under the comment ``Extended
% composition operators for function sequences'' --- plain
% composition, ``Compose head,'' and ``Compose tail'' --- each with
% its own default glyph ($\circ$, $\odot$, $\cdot$). The papers'
% own wording (``for function sequences,'' ``head,'' ``tail'') reads
% as describing composition specifically along a \emph{sequence} of
% functions, composing from one end or the other, which suggests
% these may be exactly as mathematically distinct from plain
% composition as \cs{catSeqName} is from \cs{catName} --- i.e.\ they
% may need their own macros precisely \emph{because} sequence
% composition is a different operation from pairwise composition.
% But this is not certain from the comments alone. Suggestively, the
% papers also contain a disabled (commented-out) \texttt{funcname}/
% \texttt{funcseqname} pair under the comment ``Select font for
% functions and sequences of functions'' --- the same
% single-vs.-sequence split found in the category macros, just never
% finished for functions/morphisms. That is further (if inconclusive)
% evidence that a \cs{morphSeqName} sibling to \cs{morphName}, and/or
% \cs{morphComposeHead}/\cs{morphComposeTail} siblings to
% \cs{morphCompose}, may belong in this package. None of that is
% built here; only \cs{morphCompose} is provided for now, pending the
% author's mathematical judgment rather than a guess.
% \end{quote}
%
% \subsection{Charts, atlases, and coordinates}
%
% \DescribeMacro{\chartName}
% \cs{chartName}\Arg{text} is the display name of a local coordinate
% chart.
%
% \DescribeMacro{\atlasName}
% \cs{atlasName}\Arg{text} is the display name of an atlas of charts.
%
% \DescribeMacro{\coordName}
% \cs{coordName}\Arg{text} (Coord = Coordinate) is the display name
% of a coordinate system or coordinate function.
%
% \begin{quote}
% \textbf{Placeholder, thin evidence:} the source papers use only a
% predicate \texttt{isAtl} (``Propositional function isAtl(A,E,C)'')
% and an untyped \texttt{Atl} macro (``Category of atlases from E to
% C'') for this cluster --- there is no existing inline precedent for
% chart/atlas/coordinate \emph{naming} to generalize from, unlike the
% category and morphism families above. These three commands are
% provided as minimal, unstyled hooks (default presentation is a
% no-op) so the semantic distinction exists in document markup from
% the start. Whether ``coordinate'' needs sub-kinds (e.g.\ a
% coordinate \emph{function} vs.\ a coordinate \emph{value} vs.\ a
% coordinate \emph{index}) is open and should be settled against real
% usage
% before this cluster is trusted as final.
% \end{quote}
%
% \end{documentation}
%
% \begin{implementation}
%
% \section{Implementation}
%
%<*package>
% \begin{macrocode}
\ProvidesExplPackage{semantic-markup}{2026-08-27}{0.1.0}
  {Semantic markup macros for LaTeX (LSM)}
\RequirePackage{xparse}
\RequirePackage{mathrsfs}
\RequirePackage{bm}
\RequirePackage{amssymb}
% \end{macrocode}
%
% \subsection{Generic}
%
% \begin{macro}{\lsm_sem_name:n}
% Default presentation: unstyled.
% \begin{macrocode}
\cs_new_protected:Npn \lsm_sem_name:n #1 { #1 }
\NewDocumentCommand \semName { m } { \lsm_sem_name:n {#1} }
% \end{macrocode}
% \end{macro}
%
% \subsection{Categories}
%
% \begin{macro}{\lsm_cat_name:n}
% Default presentation: script, via \cs{mathscr}.
% \begin{macrocode}
\cs_new_protected:Npn \lsm_cat_name:n #1 { \mathscr{#1} }
\NewDocumentCommand \catName { m } { \lsm_cat_name:n {#1} }
% \end{macrocode}
% \end{macro}
%
% \begin{macro}{\lsm_cat_seqname:n}
% A genuinely distinct function from \cs{lsm_cat_name:n}, not a
% call to it with a different flag: see the documentation above.
% Default presentation: bold script, via \cs{bm} and \cs{mathscr}.
% \begin{macrocode}
\cs_new_protected:Npn \lsm_cat_seqname:n #1 { \bm{\mathscr{#1}} }
\NewDocumentCommand \catSeqName { m } { \lsm_cat_seqname:n {#1} }
% \end{macrocode}
% \end{macro}
%
% \subsection{Objects}
%
% \begin{macro}{\lsm_obj_name:n}
% Default presentation: unstyled (see documentation).
% \begin{macrocode}
\cs_new_protected:Npn \lsm_obj_name:n #1 { #1 }
\NewDocumentCommand \objName { m } { \lsm_obj_name:n {#1} }
% \end{macrocode}
% \end{macro}
%
% \subsection{Morphisms}
%
% \begin{macro}{\lsm_morph_name:n}
% Default presentation: unstyled.
% \begin{macrocode}
\cs_new_protected:Npn \lsm_morph_name:n #1 { #1 }
\NewDocumentCommand \morphName { m } { \lsm_morph_name:n {#1} }
% \end{macrocode}
% \end{macro}
%
% \begin{macro}{\lsm_morph_mono:}
% Default presentation: \cs{hookrightarrow} (from \pkg{amssymb}).
% \begin{macrocode}
\cs_new_protected:Npn \lsm_morph_mono: { \hookrightarrow }
\NewDocumentCommand \morphMono { } { \lsm_morph_mono: }
% \end{macrocode}
% \end{macro}
%
% \begin{macro}{\lsm_morph_epi:}
% Default presentation: \cs{twoheadrightarrow} (from \pkg{amssymb}).
% Kept as a distinct function from \cs{lsm_morph_mono:}: see the
% documentation above.
% \begin{macrocode}
\cs_new_protected:Npn \lsm_morph_epi: { \twoheadrightarrow }
\NewDocumentCommand \morphEpi { } { \lsm_morph_epi: }
% \end{macrocode}
% \end{macro}
%
% \begin{macro}{\lsm_morph_compose:n}
% Default presentation: $\circ$, or $\circ$ with an overset label
% when a label is supplied. See the open question in the
% documentation before adding head/tail variants.
% \begin{macrocode}
\cs_new_protected:Npn \lsm_morph_compose:n #1
  {
    \tl_if_empty:nTF {#1}
      { \mathbin{\circ} }
      { \mathbin{\overset{#1}{\circ}} }
  }
\NewDocumentCommand \morphCompose { O{} } { \lsm_morph_compose:n {#1} }
% \end{macrocode}
% \end{macro}
%
% \subsection{Charts, atlases, and coordinates}
%
% \begin{macro}{\lsm_chart_name:n}
% Placeholder presentation: unstyled. See the documentation above.
% \begin{macrocode}
\cs_new_protected:Npn \lsm_chart_name:n #1 { #1 }
\NewDocumentCommand \chartName { m } { \lsm_chart_name:n {#1} }
% \end{macrocode}
% \end{macro}
%
% \begin{macro}{\lsm_atlas_name:n}
% Placeholder presentation: unstyled. See the documentation above.
% \begin{macrocode}
\cs_new_protected:Npn \lsm_atlas_name:n #1 { #1 }
\NewDocumentCommand \atlasName { m } { \lsm_atlas_name:n {#1} }
% \end{macrocode}
% \end{macro}
%
% \begin{macro}{\lsm_coord_name:n}
% Placeholder presentation: unstyled. See the documentation above.
% \begin{macrocode}
\cs_new_protected:Npn \lsm_coord_name:n #1 { #1 }
\NewDocumentCommand \coordName { m } { \lsm_coord_name:n {#1} }
% \end{macrocode}
% \end{macro}
%
%</package>
% \end{implementation}
%
% \PrintIndex
